\section{Formalization Obstacles as Mathematical Lemmas}
\label{sec:obstacles}

The development repeatedly encountered situations where a sentence that is
acceptable in a paper proof was not a well-formed theorem in Lean.  The useful
lesson is not that Lean is fussy.  Rather, the formalization identifies the
mathematical lemma that the sentence was abbreviating.  This section collects
five such obstacles.  They are the main reason the final theorem is stronger
than a certificate-consuming assembly theorem: each obstacle was eventually
turned into a constructor or theorem generated from compact support,
chartwise smoothness, and half-space geometry.

\begin{table}[t]
\centering
\scriptsize
\setlength{\tabcolsep}{3pt}
\begin{tabular}{p{0.19\textwidth}p{0.25\textwidth}p{0.24\textwidth}p{0.15\textwidth}}
\hline
Paper phrase & Ill-shaped Lean theorem & Corrected theorem/API & Role in Stokes \\
\hline
``smooth on the boundary box'' &
ambient smoothness on an open neighborhood of a closed box touching
\(x_0=0\) &
interior Stokes fields and open-interior/closed-box regularity &
legal boundary local Stokes \\

``the chart representative is supported in the box'' &
support of the totalized chart representative &
zero-localized representative and selected zero-support theorem &
support theorem for the correct representative \\

``the artificial faces vanish'' &
face integrals are zero by informal inspection &
face-coefficient support predicate and artificial-face remainder theorems &
removes all non-boundary box faces \\

``choose finitely many charts'' &
an arbitrary finite cover disconnected from partition support &
open shrinks, selected finite covers, and support-controlled partitions &
turns compact support into usable local data \\

``sum the local identities'' &
local sums and endpoint definitions use different finite indices &
generated represented endpoints over the same refined pieces &
global represented equality \\
\hline
\end{tabular}
\caption{Five paper-level phrases and the precise theorem shapes needed in
Lean.  The final first-principles theorem is small because the corrected
theorems are generated internally rather than exposed as public hypotheses.}
\label{tab:obstacle-map}
\end{table}

\Cref{tab:obstacle-map} is the organizing principle for this section.  Each
row is an instance of the same pattern.  The informal phrase is not false in
ordinary mathematical practice; it is a compressed reference to a more precise
lemma.  A certificate-consuming formalization can stop by asking the user to
provide that lemma.  The contribution of the final route is that the lemma is
proved from the preceding data and no longer appears as a theorem-facing
field.

\subsection{Boundary smoothness is not ambient closed-box smoothness}

The informal local proof near a boundary point often says that a coordinate
representative is smooth on a box.  If the model is a half-space, the word
``box'' hides a choice.  The closed coordinate box may touch the boundary
face \(x_0=0\).  Requiring an ambient open neighborhood of the closed box
inside the chart target is then too strong: it asks for Euclidean data across
the boundary, where the boundary chart is not intended to provide it.

The corrected mathematical lemma separates the roles of the box:
\begin{itemize}
  \item differentiability for the bulk term is needed on the open interior
  \(\prod_i(a_i,b_i)\);
  \item continuity and integrability are needed on the closed box and its
  faces;
  \item the lower normal face \(x_0=0\) is allowed to lie on the model
  boundary.
\end{itemize}
The Lean record \code{HalfSpaceBoxInteriorStokesFields} is exactly this
separation.  It is not a bureaucratic certificate; it is the local theorem
shape that avoids an invalid boundary-chart smoothness hypothesis.

The corrected theorem is not weaker; it is better aligned with the analysis.
It asks for differentiability exactly where the bulk derivative is evaluated
and for closed-box regularity exactly where the face integrals and
measure-theoretic bridges need it.  This is the difference between assuming a
Euclidean extension across the boundary and proving a half-space theorem.

\subsection{Support in a chart is not support of a total function}

The phrase ``the chart representative is supported in the box'' normally means
that the representative is supported in the box where the chart expression is
meaningful.  Lean functions are total, so the ordinary representative has
values outside the meaningful chart region.  If a theorem quantifies over the
ordinary support of that total function, it may be asking for facts about
values introduced only to make the function total.

The corrected lemma changes the representative.  This is stronger than adding
an auxiliary hypothesis to the old theorem.  The old theorem speaks about a
total function whose off-source values are artifacts of totalization; the new
theorem speaks about the local representative intended by the proof.  The
development defines a zero-localized representative which agrees with the
ordinary representative on the chart source and is zero outside.  The support
theorem is then a genuine topological statement:
\[
  \tsupp(\operatorname{ZeroRep}_{x,x}(\rho\omega))
  \subseteq Q(a,b).
\]
The Lean theorem
\code{zeroSupportOfSelectedSmoothRefinement} proves this statement for the
actual coefficients generated by the selected smooth refinement.

This is one of the main places where the formalization discovers a better
mathematical interface.  Asking the user for a support assumption on
\code{transitionPullbackInChart} would have hidden the problem.  Proving
\code{zeroSupportOfSelectedSmoothRefinement} removes the assumption by proving
that the generated representative has the support property that the informal
proof silently uses.

The theorem therefore has two parts that must travel together.  First, the
zero-localized representative agrees with the ordinary representative on the
source where the local integral is evaluated.  Second, it has global support
properties because it is zero outside the source.  Without the agreement
theorem the proof would be about the wrong local form; without the support
theorem artificial-face cancellation would still be a public assumption.

\subsection{Artificial faces vanish by support, not by notation}

When Stokes is applied on a coordinate box, every face of the box appears.  In
a boundary chart only the lower normal face is the manifold boundary.  The
other faces are artifacts of the localization.  A paper proof may say that
they vanish because the support is inside the box.  Lean requires the exact
disjointness statement for each face coefficient and each artificial face.

The corrected lemma has the following structure:
\[
  \tsupp(\alpha)\subseteq Q(a,b)
  \Longrightarrow
  \text{every artificial face integral of }\alpha\text{ is }0.
\]
This is implemented through the face-support API in the half-space layer.  The
global proof does not assume artificial-face cancellation; it derives it from
the zero-localized support theorem for each refined local representative.

The important point is that the support statement is made for face
coefficients, not only for the form as an abstract object.  The box theorem
integrates coefficients on coordinate faces.  The formalization proves that
evaluating the form on a fixed face frame does not enlarge topological
support, then uses the half-space support region to show that the artificial
face set is disjoint from those coefficient supports.

\subsection{Compactness is a construction principle}

Compact support is not merely a hypothesis that allows one to write a finite
sum.  The proof uses it to construct the finite index type, the selected
charts, the open shrinks, the partition coefficients, and the boundary
refinement.  If any of these choices is made independently, later support and
endpoint reconstruction statements can fail to line up.

The corrected lemma is a compactness-to-data theorem: compactness of \(K\)
together with pointwise chart-box neighborhoods produces a finite selected
cover with a support-controlled partition.
The selected cover remembers both the open sets used for the partition and the
closed or half-space boxes used for support.  This dual memory is what lets the
local support facts feed into half-space Stokes.

Thus compactness is used to generate an indexed object with invariants, not
only to prove existence of finitely many open sets.  The finite selected data
must remember enough information for three later consumers: the partition of
unity, the half-space local theorem, and the endpoint summation.

\subsection{Represented endpoints are finite coordinate theorems}

The last hidden step is endpoint reconstruction.  Once each local piece
satisfies Stokes, the proof still has to show that the finite sum of local
bulk terms is the generated represented bulk endpoint, and similarly for the
boundary terms.  This is not a definition-free conclusion: the local pieces
are sigma-indexed by boundary chart and refined box data.

The corrected theorem keeps the endpoint generated from the same finite data
used by the local proof.  This prevents a common mismatch: proving local
equalities over one index family and defining the global endpoint over a
slightly different one.  In the final route, the represented endpoints are
canonical because they are induced by the selected finite route itself.

This is why the represented endpoint layer is a mathematical layer rather
than a naming layer.  The endpoint is not chosen after the local proof; it is
the finite sum whose index family is shared by the cover, the refinement, the
zero-support theorem, and the local Stokes theorem.

\subsection{Summary of the compression}

Each obstacle has the same final form:
\[
  \text{informal paper step}
  \quad\leadsto\quad
  \text{precise generated Lean theorem}.
\]
The final first-principles API is small because these generated theorems
replace the large certificates that earlier routes exposed.  This is the
central mathematical engineering pattern of the project.
